\documentclass[11pt, a4paper]{article}

% bfw-style.tex — gemeinsame Style-Definitionen für alle Handouts/Aufgabenhefte
% Voraussetzung: Kompilation mit xelatex (wegen fontspec)

\usepackage[ngerman]{babel}
% Babel-ngerman macht " zu einem aktiven Shorthand (z.B. für "a -> ä). In unseren
% Texten verwenden wir " als normales Zeichen — daher Shorthand komplett ausschalten.
\shorthandoff{"}
\usepackage{geometry}
\geometry{a4paper, top=2.4cm, bottom=2.4cm, left=2.3cm, right=2.3cm,
          headheight=15pt, headsep=0.7cm, footskip=1.1cm}

\usepackage{fontspec}
% Fallback-Kette: Source Sans Pro -> Calibri -> Segoe UI -> Latin Modern Sans
% (Latin Modern ist immer mit MiKTeX vorhanden)
\IfFontExistsTF{Source Sans Pro}{%
  \setmainfont{Source Sans Pro}[Ligatures=TeX]%
}{\IfFontExistsTF{Calibri}{%
  \setmainfont{Calibri}[Ligatures=TeX]%
}{\IfFontExistsTF{Segoe UI}{%
  \setmainfont{Segoe UI}[Ligatures=TeX]%
}{%
  \setmainfont{Latin Modern Sans}[Ligatures=TeX]%
}}}
\IfFontExistsTF{Source Code Pro}{%
  \setmonofont{Source Code Pro}[Scale=0.92]%
}{\IfFontExistsTF{Consolas}{%
  \setmonofont{Consolas}[Scale=0.92]%
}{%
  \setmonofont{Latin Modern Mono}[Scale=0.92]%
}}

% Gesperrte Versalien für Labels/Titelbalken (Kapitälchen-Ersatz, funktioniert
% mit jeder OpenType-Schrift — Latin Modern Sans hat keine echten Kapitälchen)
\newcommand{\bfwlabelfont}{\footnotesize\bfseries\addfontfeatures{LetterSpace=8.0}}

\usepackage{xcolor}
% Farbpalette: hell, freundlich, modern
\definecolor{bfwPrimary}{HTML}{0EA5A4}    % Petrol/Türkis - Primär
\definecolor{bfwPrimaryDark}{HTML}{0F766E} % dunkler Petrol für Headings
\definecolor{bfwAccent}{HTML}{F59E0B}     % Amber - Akzent
\definecolor{bfwSuccess}{HTML}{10B981}    % Grün - Erfolg / Lösung
\definecolor{bfwWarn}{HTML}{F97316}       % Orange - Achtung
\definecolor{bfwInk}{HTML}{0F172A}        % fast Schwarz
\definecolor{bfwInkSoft}{HTML}{475569}    % gedämpftes Grau
\definecolor{bfwBgSoft}{HTML}{F8FAFC}     % off-white
\definecolor{bfwBgTeal}{HTML}{ECFEFF}     % helles Türkis
\definecolor{bfwBgAmber}{HTML}{FEF3C7}    % helles Gelb
\definecolor{bfwBgRose}{HTML}{FFE4E6}     % helles Rosé für Eselsbrücken
\definecolor{bfwTabHead}{HTML}{E4F3F2}    % zarte Kopfzeilen-Hinterlegung für Tabellen

\usepackage[colorlinks=true, linkcolor=bfwPrimaryDark, urlcolor=bfwPrimaryDark]{hyperref}
\usepackage{lastpage}
\usepackage{enumitem}
\setlist[itemize]{leftmargin=*, itemsep=2pt, topsep=2pt}
\setlist[enumerate]{leftmargin=*, itemsep=2pt, topsep=2pt}

\usepackage[most]{tcolorbox}
\usepackage{booktabs}
\usepackage{colortbl}
\usepackage{tikz}
\usetikzlibrary{decorations.pathmorphing, positioning, shapes.geometric, arrows.meta}

% ---------------------------------------------------------------------------
% Einheitliches Tabellen-Design
%   - etwas mehr Zeilenluft, dezente graue Linien (wirkt auch mit \hline ruhiger)
%   - Kopfzeile einer Tabelle bei Bedarf zart hinterlegen: \rowcolor{bfwTabHead}
% ---------------------------------------------------------------------------
\renewcommand{\arraystretch}{1.25}
\arrayrulecolor{bfwInkSoft!50}
\setlength{\heavyrulewidth}{1pt}
\setlength{\lightrulewidth}{0.5pt}

% ---------------------------------------------------------------------------
% Boxen — klare farbige Titelbalken mit gesperrten Versal-Labels
% (Titel bleibt per Option überschreibbar: \begin{merksatz}[title={...}])
% ---------------------------------------------------------------------------
\tcbset{bfwbox/.style={%
  enhanced, breakable,
  boxrule=0.5pt, arc=2.5pt,
  left=10pt, right=10pt, top=8pt, bottom=8pt,
  toptitle=4pt, bottomtitle=4pt,
  titlerule=0pt,
  fonttitle=\bfwlabelfont,
  coltitle=white
}}

% Merkkästchen — wichtige Definition / Kernpunkt
\newtcolorbox{merksatz}[1][]{%
  bfwbox,
  colback=bfwBgTeal,
  colframe=bfwPrimaryDark,
  colbacktitle=bfwPrimaryDark,
  title={MERKE},
  #1
}

% Eselsbrücke — Mnemoniks
\newtcolorbox{eselsbruecke}[1][]{%
  bfwbox,
  colback=bfwBgRose,
  colframe=bfwAccent!85!black,
  colbacktitle=bfwAccent!85!black,
  title={ESELSBRÜCKE},
  #1
}

% Beispiel-Box
\newtcolorbox{beispiel}[1][]{%
  bfwbox,
  colback=bfwBgSoft,
  colframe=bfwInkSoft,
  colbacktitle=bfwInkSoft,
  title={BEISPIEL},
  #1
}

% Lösungs-Box (für Aufgabenhefte)
\newtcolorbox{loesung}[1][]{%
  bfwbox,
  colback=bfwSuccess!7,
  colframe=bfwSuccess!65!black,
  colbacktitle=bfwSuccess!65!black,
  title={LÖSUNG},
  #1
}

% Achtung-Box — typische Fehler / Stolperfallen
\newtcolorbox{achtung}[1][]{%
  bfwbox,
  colback=bfwWarn!6,
  colframe=bfwWarn!85!black,
  colbacktitle=bfwWarn!85!black,
  title={ACHTUNG},
  #1
}

% Formel-Box — für zentrale Formeln (groß, ruhig, ohne Titelbalken)
\newtcolorbox{formelbox}[1][]{%
  enhanced,
  colback=bfwBgTeal,
  colframe=bfwPrimaryDark,
  boxrule=1pt, arc=3pt,
  left=10pt, right=10pt, top=6pt, bottom=6pt,
  #1
}

% Glossar-Eintrag
\newcommand{\glossareintrag}[2]{%
  \par\noindent\textbf{\color{bfwPrimaryDark}#1} \enspace #2\par\smallskip
}

% ---------------------------------------------------------------------------
% Abschnitts-Design: farbiges Nummern-Quadrat + feine Linie unter \section,
% dezent farbige \subsection/\subsubsection
% ---------------------------------------------------------------------------
\usepackage{titlesec}
\newcommand{\bfwSecBadge}[1]{%
  \begin{tikzpicture}[baseline={([yshift=-0.62em]n.center)}]
    \node[fill=bfwPrimaryDark, text=white, font=\normalsize\bfseries,
          minimum width=1.55em, minimum height=1.55em, inner sep=1pt,
          rounded corners=1.5pt] (n) {#1};
  \end{tikzpicture}}
\titleformat{\section}
  {\Large\bfseries\color{bfwPrimaryDark}}
  {\bfwSecBadge{\thesection}}{0.6em}{}
  [\vspace{-0.2em}{\color{bfwPrimary!60}\titlerule[0.8pt]}]
\titleformat{\subsection}
  {\large\bfseries\color{bfwPrimaryDark}}
  {\textcolor{bfwPrimary}{\thesubsection}}{0.5em}{}
\titleformat{\subsubsection}
  {\normalsize\bfseries\color{bfwInk}}
  {\textcolor{bfwPrimary}{\thesubsubsection}}{0.4em}{}
\titlespacing*{\section}{0pt}{1.6em}{1em}
\titlespacing*{\subsection}{0pt}{1.2em}{0.5em}

% TikZ-Stil "skizzenhaft" — etwas weicher als Rough.js, aber stabil
\tikzset{
  roughbox/.style = {
    draw=bfwPrimaryDark, thick, fill=bfwBgTeal,
    rounded corners=4pt, inner sep=6pt
  },
  roughaccent/.style = {
    draw=bfwAccent, thick, fill=bfwBgAmber,
    rounded corners=4pt, inner sep=6pt
  },
  roughline/.style = {
    draw=bfwInkSoft, thick, -{Stealth[length=2.5mm]}
  }
}

% Bullet-Symbole (bewusst kein FontAwesome — vermeidet Font-Installations-Probleme)
\newcommand{\faLightbulb}{$\bullet$}
\newcommand{\faBookmark}{$\bullet$}

% ---------------------------------------------------------------------------
% Kopf-/Fußzeile
%   Kopf links:  Wortlogo „Wirtschaftsfachwirt"
%   Kopf rechts: Dokument-Kurztext, gesetzt via \kopfzeile{...}
%   Fuß links:   © Can Siebert · 2026 — Fuß rechts: Seite X von Y
%   Titelseite (Seite 1) bleibt ohne Kopf-/Fußzeile (\handouttitel setzt
%   \thispagestyle{empty}).
% ---------------------------------------------------------------------------
\usepackage{fancyhdr}
\newcommand{\bfwKopfzeileText}{}
\newcommand{\kopfzeile}[1]{\renewcommand{\bfwKopfzeileText}{#1}}
\pagestyle{fancy}
\fancyhf{}
\fancyhead[L]{\small\bfseries\color{bfwPrimaryDark}Wirtschaftsfachwirt}
\fancyhead[R]{\small\color{bfwInkSoft}\bfwKopfzeileText}
\renewcommand{\headrulewidth}{0.6pt}
\renewcommand{\headrule}{{\color{bfwPrimary}\hrule width\headwidth height 0.6pt}}
\fancyfoot[L]{\small\color{bfwInkSoft}\copyright\ Can Siebert\,\textperiodcentered\,2026}
\fancyfoot[R]{\small\color{bfwInkSoft}Seite \thepage\ von \pageref*{LastPage}}
% Auch \thispagestyle{plain} (z.B. durch \maketitle) soll leer bleiben:
\fancypagestyle{plain}{\fancyhf{}\renewcommand{\headrulewidth}{0pt}\renewcommand{\headrule}{}}

% ---------------------------------------------------------------------------
% \handouttitel{Obertitel}{Untertitel}{Kurs-Label}{Termin-Zeile}
%   Einheitlicher professioneller Titelblock: Dachzeile, farbige Headline,
%   Untertitel, farbige Trennlinie und dezente Meta-Zeile (Kurs/Termin/Dozent).
%   Ersetzt \title/\maketitle. Direkt nach \begin{document} aufrufen.
% ---------------------------------------------------------------------------
\newcommand{\handouttitel}[4]{%
  \thispagestyle{empty}%
  \begingroup
  \setlength{\parindent}{0pt}%
  {\bfwlabelfont\color{bfwPrimary}WIRTSCHAFTSFACHWIRT (IHK)\par}
  \vspace{0.55em}
  {\huge\bfseries\color{bfwPrimaryDark}#1\par}
  \vspace{0.5em}
  {\large\color{bfwInkSoft}#2\par}
  \vspace{0.9em}
  {\color{bfwPrimary}\rule{\linewidth}{2pt}}\par
  \vspace{0.55em}
  {\small\color{bfwInkSoft}#3\ \,\textperiodcentered\,\ #4\ \,\textperiodcentered\,\ Dozent: Can Siebert\par}
  \endgroup
  \vspace{1.4em}%
}


\title{\vspace*{-1.5cm}\Huge\bfseries\color{bfwPrimaryDark}Aufgabenheft Lektion~2\\[0.4em]
  \large\color{bfwInkSoft}\textnormal{Wettbewerbspolitik, Staatseingriffe \& VGR --- Übungen im IHK-Klausurformat}}
\author{\textsf{Wirtschaftsfachwirt --- 1.1 Volkswirtschaftliche Grundlagen \textperiodcentered{} \small\copyright\ Can Siebert\,\textperiodcentered\,2026}}
\date{Selbstlernmaterial}

\begin{document}
\maketitle
\thispagestyle{empty}

\begin{merksatz}[title={\bfseries So nutzen Sie dieses Aufgabenheft}]
Bearbeiten Sie alle Aufgaben zunächst \emph{ohne} in die Lösungen zu schauen. Schreiben
Sie Ihre Antworten in vollständigen Sätzen --- die IHK-Prüfung verlangt formulierte
Begründungen, nicht bloße Stichworte. Beachten Sie die \emph{Operatoren} (nennen,
erläutern, berechnen, beurteilen, begründen) --- sie geben den geforderten
Antwort-Tiefegrad vor. Zeichnen Sie bei den Diagramm-Aufgaben tatsächlich auf Papier ---
in der Prüfung müssen Sie Preis-Mengen-Diagramme selbst skizzieren und beschriften. Erst
danach öffnen Sie den Lösungsteil ab Seite~\pageref{loesungen}. Ergänzend: Übungsband
Volkswirtschaft, Kapitel~1--3.
\end{merksatz}

\tableofcontents
\vspace{1em}
\hrule

\section{Aufgaben}

\subsection*{Aufgabe~1 --- Funktionen des Wettbewerbs (10 Punkte)}

Der Geschäftsführer eines mittelständischen Maschinenbauers klagt: „Der Wettbewerb setzt
uns ständig unter Druck --- ohne ihn wäre alles einfacher.``

\begin{enumerate}[label=(\alph*)]
  \item \textbf{Definieren} Sie den Begriff Wettbewerb. \hfill (2~P.)
  \item \textbf{Nennen und erläutern} Sie vier Funktionen des Wettbewerbs. \hfill (8~P.)
\end{enumerate}

\subsection*{Aufgabe~2 --- Kartellverbot und Kartellarten (12 Punkte)}

Vier Zementhersteller vereinbaren mündlich, (1) ihre Preise künftig einheitlich zu
erhöhen, (2) das Bundesgebiet regional unter sich aufzuteilen und (3) die
Produktionsmengen zu drosseln, um das Angebot zu verknappen.

\begin{enumerate}[label=(\alph*)]
  \item \textbf{Definieren} Sie den Begriff Kartell und \textbf{erläutern} Sie, was mit den beteiligten Unternehmen in rechtlicher und wirtschaftlicher Hinsicht geschieht. \hfill (4~P.)
  \item \textbf{Ordnen} Sie die drei Vereinbarungen jeweils einer Kartellart zu. \hfill (3~P.)
  \item \textbf{Erläutern} Sie unter Angabe der Rechtsgrundlage, warum diese Vereinbarungen verboten sind, obwohl sie nur mündlich getroffen wurden. \hfill (3~P.)
  \item \textbf{Nennen} Sie zwei Gruppen von Ausnahmen vom Kartellverbot. \hfill (2~P.)
\end{enumerate}

\subsection*{Aufgabe~3 --- Missbrauchsaufsicht und Fusionskontrolle (12 Punkte)}

\begin{enumerate}[label=(\alph*)]
  \item \textbf{Nennen} Sie die drei Säulen des GWB und die Behörde, die über die Einhaltung des GWB wacht. \hfill (4~P.)
  \item Ein marktbeherrschender Handelskonzern fordert von einem wirtschaftlich abhängigen Zulieferer extreme Preisnachlässe; ein anderes marktbeherrschendes Unternehmen verkauft gezielt zu Kampfpreisen, um einen neuen Konkurrenten vom Markt zu verdrängen. \textbf{Ordnen} Sie beide Fälle der zutreffenden Missbrauchsform zu und \textbf{begründen} Sie. \hfill (4~P.)
  \item \textbf{Erläutern} Sie, warum die Fusionskontrolle als vorbeugende Maßnahme gilt, und \textbf{erklären} Sie den Begriff Ministererlaubnis. \hfill (4~P.)
\end{enumerate}

\subsection*{Aufgabe~4 --- Mindestpreis (12 Punkte)}

Der Staat legt für ein Agrarprodukt einen Mindestpreis fest, der über dem
Gleichgewichtspreis liegt.

\begin{enumerate}[label=(\alph*)]
  \item \textbf{Skizzieren} Sie die Situation in einem Preis-Mengen-Diagramm (Achsen, Angebots- und Nachfragekurve, Gleichgewichtspreis, Mindestpreis, Mengen beschriften). \hfill (5~P.)
  \item \textbf{Erläutern} Sie, welche Marktseite geschützt werden soll und welche Marktsituation durch den Mindestpreis entsteht. \hfill (4~P.)
  \item \textbf{Beschreiben} Sie am Beispiel des EU-Agrarmarktes, welche Folgeprobleme und Folgemaßnahmen sich für den Staat ergeben. \hfill (3~P.)
\end{enumerate}

\subsection*{Aufgabe~5 --- Höchstpreis und Beurteilung (12 Punkte)}

Um Mieter zu schützen, wird auf einem angespannten Wohnungsmarkt eine Mietobergrenze
unterhalb der Gleichgewichtsmiete festgelegt.

\begin{enumerate}[label=(\alph*)]
  \item \textbf{Erklären} Sie den Begriff Höchstpreis und \textbf{begründen} Sie, warum er unterhalb des Gleichgewichtspreises liegen muss, um zu wirken. \hfill (3~P.)
  \item \textbf{Erläutern} Sie die Marktsituation, die durch die Mietobergrenze entsteht, und zwei typische Folgeerscheinungen. \hfill (5~P.)
  \item \textbf{Beurteilen} Sie Höchst- und Mindestpreise aus ordnungspolitischer Sicht (Stichwort: Funktionen des Marktpreises). \hfill (4~P.)
\end{enumerate}

\subsection*{Aufgabe~6 --- Marktkonforme Eingriffe (10 Punkte)}

\begin{enumerate}[label=(\alph*)]
  \item \textbf{Unterscheiden} Sie marktkonforme und marktkonträre Staatseingriffe mit je einem Beispiel. \hfill (4~P.)
  \item \textbf{Erläutern} Sie, wie sich eine Erhöhung der Tabaksteuer auf Gleichgewichtspreis, Gleichgewichtsmenge und Steuereinnahmen auswirkt und welches politische Ziel damit verfolgt werden kann. \hfill (4~P.)
  \item \textbf{Erklären} Sie, warum Dauersubventionen ordnungspolitisch problematisch sind. \hfill (2~P.)
\end{enumerate}

\subsection*{Aufgabe~7 --- VGR: Volkseinkommen, Lohn- und Gewinnquote (16 Punkte)}

Für eine Volkswirtschaft liegen für das abgelaufene Jahr folgende Daten vor (in Mrd.~€):

\begin{center}
\begin{tabular}{l r}
\toprule
\textbf{Größe} & \textbf{Mrd.~€}\\
\midrule
Arbeitnehmerentgelt (Inländer) & 1.750\\
Unternehmens- und Vermögenseinkommen (Inländer) & 750\\
Produktions- und Importabgaben abzüglich Subventionen & 330\\
Abschreibungen & 600\\
Saldo der Primäreinkommen aus der übrigen Welt & $+70$\\
\bottomrule
\end{tabular}
\end{center}

\begin{enumerate}[label=(\alph*)]
  \item \textbf{Berechnen} Sie das Volkseinkommen. \hfill (2~P.)
  \item \textbf{Berechnen} Sie die Lohnquote und die Gewinnquote und \textbf{interpretieren} Sie die Ergebnisse. \hfill (6~P.)
  \item \textbf{Berechnen} Sie ausgehend vom Volkseinkommen das Bruttonationaleinkommen (BNE) und daraus das Bruttoinlandsprodukt (BIP). \hfill (5~P.)
  \item \textbf{Erläutern} Sie den konzeptionellen Unterschied zwischen BIP und BNE. \hfill (3~P.)
\end{enumerate}

\subsection*{Aufgabe~8 --- Nominales/reales BIP und Grenzen des BIP (12 Punkte)}

Das nominale BIP eines Landes ist gegenüber dem Vorjahr um 3,4~\% gestiegen; die
Preissteigerungsrate betrug im gleichen Zeitraum 1,9~\%.

\begin{enumerate}[label=(\alph*)]
  \item \textbf{Unterscheiden} Sie nominales und reales BIP. \hfill (4~P.)
  \item \textbf{Ermitteln} Sie näherungsweise das reale Wirtschaftswachstum und \textbf{begründen} Sie, warum für Wachstumsaussagen das reale BIP herangezogen wird. \hfill (4~P.)
  \item \textbf{Nennen} Sie vier Kritikpunkte an der Aussagekraft des BIP als Wohlstandsindikator. \hfill (4~P.)
\end{enumerate}

\subsection*{Aufgabe~9 --- Magisches Viereck und Verteilungspolitik (14 Punkte)}

\begin{enumerate}[label=(\alph*)]
  \item \textbf{Nennen} Sie die vier Ziele des magischen Vierecks nach dem Stabilitätsgesetz mit ihrer jeweiligen Messgröße. \hfill (8~P.)
  \item \textbf{Erläutern} Sie, warum das Viereck als „magisch`` bezeichnet wird, und \textbf{erklären} Sie je einen Zielkonflikt und eine Zielharmonie an einem Beispiel. \hfill (4~P.)
  \item \textbf{Nennen} Sie die beiden zusätzlichen Ziele des magischen Sechsecks. \hfill (2~P.)
\end{enumerate}

\newpage
\section{Lösungen}\label{loesungen}

\begin{loesung}[title={Lösung Aufgabe~1 --- Funktionen des Wettbewerbs}]
\textbf{(a)} Wettbewerb ist die Rivalität der Marktteilnehmer unter- und gegeneinander
mit dem Ziel, die eigenen Ziele zulasten der Konkurrenten durchzusetzen.

\textbf{(b)} (je zwei Punkte, vier aus sechs) \emph{Freiheitsfunktion:} Jeder kann seine
Fähigkeiten und Mittel für frei gewählte Zwecke einsetzen (Unternehmertätigkeit,
Konsumwahl). \emph{Kontrollfunktion:} Anbieter und Nachfrager kontrollieren sich
gegenseitig; die wirtschaftliche Macht des Einzelnen wird begrenzt.
\emph{Steuerungsfunktion (Lenkungsfunktion):} beschleunigt die Koordination von Angebot
und Nachfrage (Koordinations- und Anpassungsfunktion) und sorgt für den produktivsten
Einsatz der Produktionsfaktoren (Allokationsfunktion). \emph{Innovationsfunktion:}
Anreiz, Produkte zu verbessern und kostengünstige Verfahren sowie technischen Fortschritt
zu nutzen. \emph{Auslesefunktion:} Auslese der Marktteilnehmer nach wirtschaftlicher
Leistungsfähigkeit. \emph{Verteilungsfunktion:} Entlohnung der Produktionsfaktoren
entsprechend ihrer Leistung (leistungsorientierte Einkommensverteilung).
\end{loesung}

\begin{loesung}[title={Lösung Aufgabe~2 --- Kartellverbot und Kartellarten}]
\textbf{(a)} Kartelle sind Kooperationen von Unternehmen auf Basis vertraglicher
(auch mündlicher) Vereinbarungen zur Verbesserung der Wettbewerbsposition. Die
beteiligten Unternehmen bleiben \emph{rechtlich selbstständig}, ihre
\emph{wirtschaftliche} Selbstständigkeit wird durch den Inhalt der Vereinbarungen
teilweise eingeschränkt.

\textbf{(b)} (1) Preiskartell --- Absprache einer einheitlichen Preispolitik.
(2) Gebietskartell --- Aufteilung des Absatzmarktes. (3) Quotenkartell --- künstliche
Verknappung des Angebots durch abgesprochene Produktionsquoten.

\textbf{(c)} Nach \S\,1 GWB (europarechtlich: Art.~101 AEUV) sind Vereinbarungen und
\emph{aufeinander abgestimmte Verhaltensweisen} verboten, die eine Verhinderung,
Einschränkung oder Verfälschung des Wettbewerbs bewirken. Auf die Form kommt es nicht an
--- auch mündliche Absprachen und bloß abgestimmtes Verhalten sind erfasst. Bei Verstößen
drohen Bußgelder und Gewinnabschöpfung.

\textbf{(d)} Zwei aus drei: besondere Ausnahmeregelungen nach dem GWB (z.\,B.
Mittelstandskartelle nach \S\,3 GWB), EU-Gruppenfreistellungsverordnungen,
Einzelfreistellungen (Legalausnahmen nach \S\,2 GWB).
\end{loesung}

\begin{loesung}[title={Lösung Aufgabe~3 --- Missbrauchsaufsicht und Fusionskontrolle}]
\textbf{(a)} Kartellverbot (\S\,1 GWB), Missbrauchsaufsicht (\S\S\,19\,ff. GWB),
Fusionskontrolle (\S\,35 GWB); über die Einhaltung wacht das \emph{Bundeskartellamt}.

\textbf{(b)} Fall 1: \emph{Ausbeutungsmissbrauch} --- die Marktmacht wird genutzt, um von
einem abhängigen Zulieferer extreme Preisnachlässe zu erzwingen (bzw. von Kunden
überhöhte Preise zu fordern). Fall 2: \emph{Behinderungsmissbrauch} --- die Marktmacht
wird gezielt eingesetzt, um Konkurrenten über Kampfpreise (Dumping) zu verdrängen oder
ihnen den Marktzugang zu versperren.

\textbf{(c)} Die Missbrauchsaufsicht betrifft \emph{bestehende} marktbeherrschende
Unternehmen; die Fusionskontrolle setzt früher an: Geplante Fusionen müssen dem
Bundeskartellamt angezeigt werden, das den Zusammenschluss untersagen kann, \emph{bevor}
eine marktbeherrschende Stellung entsteht (vorbeugend). \emph{Ministererlaubnis}
(\S\,42 GWB): Der Bundeswirtschaftsminister kann einen vom Bundeskartellamt untersagten
Zusammenschluss dennoch erlauben, wenn er von gesamtwirtschaftlichem Vorteil ist und die
marktwirtschaftliche Ordnung nicht gefährdet wird.
\end{loesung}

\begin{loesung}[title={Lösung Aufgabe~4 --- Mindestpreis}]
\textbf{(a)} Skizze: x-Achse Menge (x), y-Achse Preis (p); fallende Nachfragekurve N,
steigende Angebotskurve A, Schnittpunkt beim Gleichgewichtspreis $p_0$ mit Menge $x_0$.
Der Mindestpreis $\bar{p}$ liegt als waagerechte Linie \emph{oberhalb} von $p_0$. Beim
Mindestpreis gilt: nachgefragte Menge $x_N <$ Gleichgewichtsmenge $x_0 <$ angebotene
Menge $x_A$.

\textbf{(b)} Ziel ist der \emph{Schutz der Anbieter} (z.\,B. Landwirte), die einen
auskömmlichen Preis erhalten sollen. Da der Preis über dem Gleichgewicht liegt, weiten
die Anbieter ihr Angebot aus, während die Nachfrage sinkt --- es entsteht ein
\emph{Angebotsüberschuss} (Angebotsüberhang), der zum Mindestpreis keinen Käufer findet.

\textbf{(c)} Im EU-Agrarmarkt führten Mindestpreise mit Abnahmegarantien zu deutlicher
Überproduktion. Der Staat musste die Überschüsse durch Einlagerung, Pflege, Verarbeitung,
Veräußerung oder Vernichtung beseitigen (erhebliche Kosten) und ergriff Folgemaßnahmen
zur Angebotseindämmung: Produktionskontingente und Prämien für freiwillige
Angebotsreduzierung (z.\,B. Flächenstilllegung). Bei Kontingenten kann zudem ein
\emph{grauer Markt} entstehen, auf dem die Mehrproduktion unter dem Mindestpreis
angeboten wird.
\end{loesung}

\begin{loesung}[title={Lösung Aufgabe~5 --- Höchstpreis und Beurteilung}]
\textbf{(a)} Ein Höchstpreis ist eine staatlich verordnete Preisobergrenze für ein
bestimmtes Gut. Läge er über dem Gleichgewichtspreis, würde sich weiterhin der (dann
zulässige) niedrigere Gleichgewichtspreis bilden --- der Eingriff wäre wirkungslos. Nur
unterhalb des Gleichgewichtspreises verändert er das Marktergebnis.

\textbf{(b)} Zum künstlich niedrigen Preis steigt die nachgefragte Menge, während das
Angebot zurückgeht (Vermietung lohnt sich weniger) --- es entsteht ein
\emph{Nachfrageüberschuss} (Wohnungsmangel). Typische Folgen: Es sind flankierende
Maßnahmen nötig (Subventionen, damit überhaupt Angebot besteht; Kontingentierung und
Rationierung der knappen Menge); häufig bilden sich \emph{Schwarzmärkte}, auf denen zu
höheren Preisen verkauft wird. In seltenen Fällen können Importe die Lücke schließen,
wenn ausländische Anbieter kostengünstiger anbieten.

\textbf{(c)} Beurteilung: Höchst- und Mindestpreise sind \emph{marktkonträre} Eingriffe
--- sie hebeln den Preismechanismus aus und setzen die Funktionen des Marktpreises
(Signal-, Ausgleichs-, Allokations-, Selektionsfunktion) außer Kraft. Es entstehen
dauerhafte Überschüsse bzw. Knappheiten, Folgekosten für den Staat, Fehlallokation von
Produktionsfaktoren sowie graue und schwarze Märkte. Marktkonforme Instrumente sind
ordnungspolitisch vorzuziehen.
\end{loesung}

\begin{loesung}[title={Lösung Aufgabe~6 --- Marktkonforme Eingriffe}]
\textbf{(a)} \emph{Marktkonforme (indirekte) Eingriffe} beeinflussen Angebot und
Nachfrage, ohne den Preismechanismus außer Kraft zu setzen --- z.\,B. Verbrauchssteuern,
Zölle, Subventionen, Transferleistungen (Wohngeld, Kindergeld).
\emph{Marktkonträre (direkte) Eingriffe} setzen den Preismechanismus außer Kraft ---
z.\,B. staatliche Höchst- und Mindestpreise.

\textbf{(b)} Die höhere Tabaksteuer verteuert das Produkt: Der Gleichgewichtspreis
steigt, die Gleichgewichtsmenge sinkt. Durch die geringere Absatzmenge sinken tendenziell
auch die Steuereinnahmen je nach Nachfragereaktion. Gesundheitspolitisches Ziel: Das
Verbraucherverhalten wird in eine erwünschte Richtung gelenkt (weniger Raucher, weniger
Krankheitsfälle).

\textbf{(c)} Dauersubventionen behindern sämtliche Preisfunktionen: Der künstlich
niedrige Preis sendet falsche Knappheitssignale, das Ausscheiden von Anbietern mit hohen
Produktionskosten wird verhindert (Selektionsfunktion), und Produktionsfaktoren bleiben
in der Herstellung nicht (mehr) gefragter Produkte gebunden (Fehlallokation).
\end{loesung}

\begin{loesung}[title={Lösung Aufgabe~7 --- VGR: Volkseinkommen, Lohn- und Gewinnquote}]
\textbf{(a)} Volkseinkommen $=$ Arbeitnehmerentgelt $+$ Unternehmens- und
Vermögenseinkommen $= 1.750 + 750 = \textbf{2.500 Mrd.~€}$.

\textbf{(b)} Lohnquote $= \frac{1.750}{2.500} \cdot 100 = \textbf{70~\%}$;
Gewinnquote $= \frac{750}{2.500} \cdot 100 = \textbf{30~\%}$.
Interpretation: 70~\% des Volkseinkommens entfallen auf Arbeitnehmerentgelte
(unselbstständige Arbeit), 30~\% auf Unternehmens- und Vermögenseinkommen; beide Quoten
ergänzen sich zu 100~\%. Vorsicht: Zinseinkünfte von Arbeitnehmern zählen zur
Gewinnquote, Vorstandsgehälter zum Arbeitnehmerentgelt --- die Quoten sind daher nur
begrenzt als „Verteilungsgerechtigkeits-Maß`` interpretierbar.

\textbf{(c)} BNE $=$ Volkseinkommen $+$ Produktions- und Importabgaben abzüglich
Subventionen $+$ Abschreibungen $= 2.500 + 330 + 600 = \textbf{3.430 Mrd.~€}$.
BIP $=$ BNE $-$ Saldo der Primäreinkommen aus der übrigen Welt
$= 3.430 - 70 = \textbf{3.360 Mrd.~€}$.

\textbf{(d)} Das BIP folgt dem \emph{Inlandskonzept}: Es misst die Wirtschaftsleistung
innerhalb der geografischen Grenzen --- unabhängig davon, wer sie erbringt. Das BNE folgt
dem \emph{Inländerkonzept}: Es misst das Einkommen der Inländer aus dem In- und Ausland
--- unabhängig davon, wo es erwirtschaftet wurde. Rechnerisch: BIP $-$ Inlandseinkommen
von Nichtinländern $+$ Auslandseinkommen von Inländern $=$ BNE.
\end{loesung}

\begin{loesung}[title={Lösung Aufgabe~8 --- Nominales/reales BIP und Grenzen}]
\textbf{(a)} Das \emph{nominale} BIP wird in laufenden Preisen (aktuellen Marktpreisen)
ausgewiesen --- es enthält Mengen- \emph{und} Preisänderungen. Das \emph{reale} BIP wird
zu Preisen eines festgelegten Basisjahres bewertet --- Preissteigerungen sind eliminiert,
es zeigt den echten (mengenmäßigen) Leistungszuwachs.

\textbf{(b)} Reales Wachstum $\approx$ nominales Wachstum $-$ Preissteigerungsrate
$= 3{,}4~\% - 1{,}9~\% = \textbf{1{,}5~\%}$. Nur das reale BIP zeigt, ob tatsächlich mehr
Güter und Dienstleistungen produziert wurden --- ein rein preisbedingter Anstieg des
nominalen BIP wäre kein Wachstum. Im internationalen und zeitlichen Vergleich werden
daher stets reale Werte zugrunde gelegt.

\textbf{(c)} Vier Kritikpunkte (Beispiele): (1) Eigenleistungen der Haushalte, Hausarbeit
und ehrenamtliche Tätigkeiten werden nicht erfasst; (2) Schwarzarbeit
(Schattenwirtschaft) wird nur geschätzt; (3) wohlstandsmindernde Ereignisse können das
BIP erhöhen (Reparatur von Unfall- und Umweltschäden); (4) keine Aussage über die
Verteilung des Einkommens, Umweltbelastung und Lebensqualität --- ein hohes
Pro-Kopf-Einkommen schließt Armut von Bevölkerungsteilen nicht aus.
\end{loesung}

\begin{loesung}[title={Lösung Aufgabe~9 --- Magisches Viereck und Verteilungspolitik}]
\textbf{(a)} Nach \S\,1 StabG (1967): (1) \emph{Stabilität des Preisniveaus} ---
gemessen an der Inflationsrate auf Basis des Verbraucherpreisindex (Ziel: knapp unter
2~\%); (2) \emph{hoher Beschäftigungsstand} --- gemessen an der Arbeitslosenquote (als
erreicht gilt ca. 3--4~\%); (3) \emph{außenwirtschaftliches Gleichgewicht} --- gemessen
am Außenbeitrag (Ziel: positiver Außenbeitrag von 1--2~\% des nominalen BIP);
(4) \emph{stetiges und angemessenes Wirtschaftswachstum} --- gemessen an der Zuwachsrate
des realen BIP (angemessen: rund 2~\% pro Jahr).

\textbf{(b)} „Magisch``, weil es nahezu unmöglich ist, alle Ziele gleichzeitig
vollständig zu erreichen --- sie stehen teilweise in Konkurrenz. \emph{Zielkonflikt}
(Beispiel): Starkes Wirtschaftswachstum zieht i.\,d.\,R. Preiserhöhungen nach sich und
gefährdet die Preisniveaustabilität. \emph{Zielharmonie} (Beispiel): Wirtschaftswachstum
und Vollbeschäftigung --- steigt das BIP, werden mehr Produktionsfaktoren (auch Arbeit)
benötigt, die Arbeitslosigkeit sinkt.

\textbf{(c)} Erhaltung einer lebenswerten Umwelt (Umweltschutz) und gerechte
Einkommens- und Vermögensverteilung (qualitative Ziele, nicht eindeutig messbar).
\end{loesung}

\vspace{1em}
\hrule
\vspace{0.8em}
\noindent\textsf{\small Weiterüben: Übungsband Volkswirtschaft Kap.~1--3 (mit Lösungsteil) und das
interaktive Quiz dieser Lektion. Nächste Lektion: Konjunktur, Wirtschaftspolitik \& Außenwirtschaft.}

\end{document}
